\documentclass[11pt]{article}
\PassOptionsToPackage{obeyspaces}{url}
\usepackage[margin=1in]{geometry}
\usepackage[T1]{fontenc}
\usepackage{charter}
\usepackage[scaled=0.9]{helvet}
\usepackage{courier}

\usepackage{microtype}
\usepackage{booktabs}
\usepackage{tabularx}
\usepackage{longtable}
\usepackage{xcolor}
\usepackage{listings}
\usepackage{enumitem}
\usepackage{amsmath}
\usepackage[hidelinks]{hyperref}

\definecolor{codebg}{gray}{0.96}
\definecolor{lo}{RGB}{30,120,60}
\definecolor{med}{RGB}{190,120,0}
\definecolor{hi}{RGB}{180,30,30}
\newcommand{\low}{\textcolor{lo}{\textbf{low}}}
\newcommand{\medium}{\textcolor{med}{\textbf{medium}}}
\newcommand{\high}{\textcolor{hi}{\textbf{high}}}
\usepackage{url}
\usepackage{array}
\usepackage{float}
\def\UrlBreaks{\do\/\do\_\do\.\do\-\do\:}
\DeclareUrlCommand\f{\urlstyle{tt}}
\DeclareUrlCommand\cd{\urlstyle{tt}\def\UrlBreaks{\do\.\do\(\do\)\do\,\do\_\do\:\do\<\do\>\do\|\do\/}}

\lstset{
  basicstyle=\ttfamily\footnotesize,
  backgroundcolor=\color{codebg},
  frame=none, columns=fullflexible, keepspaces=true,
  breaklines=true, xleftmargin=1em, aboveskip=0.6em, belowskip=0.6em,
  keywordstyle=\bfseries, language=C, morekeywords={fn,let,mut,impl,pub,match,Some,None,Ok,self,Self,move,async}
}
\setlist{itemsep=0.2em, topsep=0.3em}

\title{Clone and Copy Discipline in \texttt{f1r3node-rust}\\[0.3em]
\large Where the copies are, which ones matter, and how to remove them without moving consensus}
\author{F1R3FLY.io Research Note\\[0.2em]\small Prepared with Claude for Meredith}
\date{20 October 2026 \quad\textbullet\quad \texttt{dev} @ \texttt{e8008c3}}

\setlength{\emergencystretch}{3em}
\begin{document}
\maketitle

\begin{abstract}
\noindent
This note audits the \texttt{dev} branch of \texttt{f1r3node-rust} for unnecessary
\texttt{clone()}, \texttt{to\_vec()}, \texttt{cloned()} and \texttt{to\_owned()} calls, with the
aim of finding the largest performance gains that leave consensus behaviour untouched.
Production code (tests and scripts excluded) contains about 3{,}800 explicit clones in roughly
268{,}000 lines. The raw count is a poor guide: block hashes, state hashes and validator keys are
already reference-counted \texttt{Bytes}, so most casper clones are pointer bumps. The expensive
copies cluster in a small number of places where a value with deep structure --- the
protobuf \texttt{Par} tree, an \texttt{Env}, a block body, a deploy-chain index, a stream buffer ---
is copied because the code transliterates a Scala idiom (persistent, structurally shared
collections) into Rust collections that copy deeply. Seven of these sites are
\emph{asymptotic} defects (quadratic in block size, binding count, COMM count or chain length).
Microbenchmarks of the exact patterns show 17--270$\times$ slowdowns at realistic sizes, and
126$\times$ (5.6\,s versus 44\,ms) for reassembling a 64\,MiB stream. The note ranks the findings
into four tiers, separates the clones that are \emph{semantically required} (random-number
splitting, cost order, canonical sorting, hash bytes) from those that are not, and proposes a
sequence of small pull requests, each with a determinism guard.
\end{abstract}

\tableofcontents
\newpage

%=====================================================================
\section{Scope and method}

\paragraph{Object of study.} The \texttt{dev} branch of \texttt{F1R3FLY-io/f1r3node-rust} at
commit \texttt{e8008c3} (merge of PR \#781, 20 October 2026). The workspace has eleven crates;
the audit covers \texttt{rholang}, \texttt{rho-pure-eval}, \texttt{rspace++}, \texttt{casper},
\texttt{block-storage}, \texttt{models}, \texttt{comm}, \texttt{shared} and \texttt{node}.

\paragraph{Census.} A script counts the clone family of calls per crate, and separates production
from test code. A file counts as test code when its path is under \texttt{tests/},
\texttt{benches/}, \texttt{scripts/}, \texttt{fuzz/}, or a test-utility module. Inline
\cd{#[cfg(test)]} items are removed by brace matching before the count.

\paragraph{Deep reading.} The four highest-traffic subsystems --- the Rholang reducer and
substitution, the tuple space (RSpace), casper validation and merging, and the model,
transport and storage layers --- were read in full along their hot paths. Each finding records
its location, the type that is copied, a cost class, whether it lies on a hot path, a proposed
change, and a risk grade. Every line reference in this note was checked against the source at
\texttt{e8008c3}.

\paragraph{Cost classes.}
\begin{description}[leftmargin=1.5em]
\item[Deep tree] a recursive structure (for example \texttt{Par}); one heap allocation per node.
\item[Large buffer] a \texttt{Vec<u8>} or \texttt{Vec<T>} whose size grows with the input.
\item[Small allocation] a 32-byte \texttt{Vec<u8>} hash, a short \texttt{String}.
\item[Pointer bump] \texttt{Arc}, \texttt{Bytes}, or an \texttt{imbl} persistent collection.
\end{description}

\paragraph{Risk grades.} A change is \low{} risk if it only changes ownership (move instead of
copy, borrow instead of copy) and no observable output can change. It is \medium{} risk if it
changes a type or signature shared across crates, or touches code on which determinism depends
without changing the computation. It is \high{} risk if it touches hash bytes, serialisation
formats, canonical sort order, or replay semantics.

\paragraph{Limitations.} The node was not built or profiled for this note. The gains below are
static estimates, supported by standalone microbenchmarks of the exact code patterns
(Section~\ref{sec:bench}). Section~\ref{sec:measure} proposes how to confirm them with the
repository's own soak and latency tooling before and after each change.

%=====================================================================
\section{Census}

\begin{table}[ht]
\centering\small
\begin{tabular}{lrrrrr}
\toprule
Crate & Prod.\ lines & \texttt{clone()} & \texttt{to\_vec()} & \texttt{cloned()} & per kLoC\\
\midrule
casper          & 83{,}707 & 1{,}372 & 70 & 125 & 16.4\\
rholang         & 83{,}750 &   880 & 54 &  12 & 10.5\\
rspace++        & 18{,}892 &   482 & 69 &  61 & 25.5\\
node            & 32{,}198 &   432 & 14 &   7 & 13.4\\
models          & 12{,}605 &   231 & 11 &   7 & 18.3\\
block-storage   &  9{,}745 &   184 & 38 &  36 & 18.9\\
comm            & 13{,}342 &   174 & 13 &  10 & 13.0\\
shared, crypto, graphz, rho-pure-eval & 11{,}908 & 62 & 34 & 4 & 5.2\\
\midrule
\textbf{Total production} & \textbf{266{,}147} & \textbf{3{,}817} & \textbf{303} & \textbf{262} & \textbf{14.3}\\
Test code & & 7{,}045 & & & \\
\bottomrule
\end{tabular}
\caption{Clone-family calls in production code at \texttt{e8008c3}.}
\end{table}

The densest files are \f{rholang/.../interpreter/reduce.rs} (230), \f{casper/.../engine/casper_launch.rs} (125),
\f{node/.../runtime/setup.rs} (116), \f{rholang/.../rho_runtime.rs} (114) and
\f{rspace++/.../replay_rspace.rs} (78). Two of these five (\texttt{casper\_launch.rs} and
\texttt{setup.rs}) run once at start-up and copy \texttt{Arc} handles; they are irrelevant to
throughput. The count therefore tells where to look, not what to fix.

%=====================================================================
\section{What a clone costs here}
\label{sec:cost}

The cost of a clone depends on the type, and the types in this codebase fall into two clear
groups.

\paragraph{Already cheap.}
\begin{itemize}
\item \texttt{BlockHash}, \texttt{StateHash}, \texttt{Validator} and casper \texttt{ByteString}
are \cd{prost::bytes::Bytes}: \f{models/build.rs} lines 55--56 configure
\cd{.bytes(".casper")} and \cd{.bytes(".routing")}. Most of the 1{,}372 casper clones are
of these types.
\item \texttt{KeyValueDagRepresentation} and \texttt{DagState} are built on \texttt{imbl}
persistent maps; cloning them is $O(1)$.
\item \texttt{DebruijnInterpreter} (\f{reduce.rs}:231--257) is about fourteen \texttt{Arc}
fields; one clone is fourteen atomic increments and no heap work.
\item The hot store's continuations are already \cd{Vec<Arc<WaitingContinuation>>}. A comment
at \f{hot_store.rs}:408--412 records that an earlier change removed a deep clone of
persistent contract bodies on every produce. That change is the precedent this note extends.
\end{itemize}

\paragraph{Expensive.}
\begin{itemize}
\item \texttt{Par} and every Rholang AST node. The \texttt{rhoapi} package has no
\cd{.bytes(...)} setting, and \texttt{Par} is nine \texttt{repeated} message vectors that
recurse with no \texttt{Box}, \texttt{Arc} or \texttt{Rc}. Every clone is a full deep copy.
\item \texttt{Env<Par>} (\f{rho-pure-eval/src/env.rs}:12) is a \cd{HashMap<i32, Par>}. A clone
deep-copies every bound process.
\item \texttt{BlockMessage}. The hash fields are \texttt{Bytes}, but the body holds
\cd{Vec<ProcessedDeploy>}, each with a \texttt{String} term and a \cd{Vec<Event>} deploy
log. A clone costs on the order of the number of event-log entries.
\item \texttt{DeployChainIndex} carries three \texttt{EventLogIndex} values of eleven hash sets
each, plus a \texttt{StateChange} of three \texttt{DashMap}s.
\item \texttt{Blake2b256Hash} is \cd{pub struct Blake2b256Hash(pub Vec<u8>)}
(\f{rspace++/.../blake2b256_hash.rs}:17). It is not \texttt{Copy}; each clone is an allocation.
It keys the RSpace event log and the merge channel maps.
\end{itemize}

\paragraph{Diagnosis.} The repository was extracted from a Scala code base. In Scala,
\texttt{env.put(x)} on an immutable map returns a new map that shares almost all of its
structure with the old one, and passing a value ``by copy'' is free. Several of the most
expensive Rust sites reproduce that interface literally: they mutate in place \emph{and} return
a full clone, or they clone a collection, append one element, and store the result.
\texttt{Env::put}, \texttt{SortedParMap::insert}, the stream reassembly buffer and
\texttt{EventLogIndex::combine} all follow this pattern. The general remedy has two forms.
Where the old value is dead after the call, move it. Where the old value is genuinely shared,
use a persistent structure (\texttt{imbl} is already a workspace dependency) or \texttt{Arc}.
The repository's own guideline in \texttt{CLAUDE.md} asks for ``zero-cost abstractions, proper
ownership''; the findings below are the places where the port has not yet met it.

%=====================================================================
\section{Findings}

The findings are in four tiers. Tier~0 holds asymptotic defects, where a fix changes the growth
rate. Tier~1 holds hot-path, low-risk, mechanical changes. Tier~2 holds structural changes of
medium risk with large gains. Tier~3 holds changes that are not recommended now.

\subsection{Tier 0: asymptotic defects}

\paragraph{0.1 Stream reassembly is quadratic in message size.}
\f{comm/src/rust/transport/stream_handler.rs}:482--491. For each incoming 256\,KiB chunk, the
whole buffer received so far is copied out of the cache, extended, and written back.
\begin{lstlisting}
let received_bytes = content_data.to_vec();
let mut existing_data = cache.get(&streamed.key)
    .map(|entry| entry.value().clone()).unwrap_or_default();
existing_data.extend_from_slice(&received_bytes);
cache.insert(streamed.key.clone(), existing_data);
\end{lstlisting}
A 10\,MiB block arrives in 40 chunks and causes about 200\,MiB of copying; a 64\,MiB stream
causes about 8\,GiB. \texttt{restore} (line~393) then clones the completed buffer again before it
deletes the entry. \emph{Fix:}
\cd{cache.entry(key).or_default().extend_from_slice(&content_data)}, and take the buffer at
restore with \texttt{cache.remove}. Path: every received block and approved block.
Risk: \low{}; internal only, wire bytes unchanged.

\paragraph{0.2 Replay scans its COMM list with a clone per step.}
\f{rspace++/src/rspace/replay_rspace.rs}:1126--1153, \texttt{get\_comm\_or\_candidate}. Each
loop iteration runs \cd{go(cs.clone())} on the remaining \cd{Vec<COMM>}, and the miss branch
returns \cd{rem.to_vec()}. A scan over $n$ COMMs therefore copies $O(n^2)$ COMMs. In
addition, \texttt{matches} and \texttt{was\_repeated\_enough\_times} (lines 966, 973) take
\texttt{COMM} and \texttt{Datum} by value, so callers clone both for every datum, every channel
and every candidate COMM (lines 808, 959). \emph{Fix:} iterate \cd{comms.iter()} and return the
first match; make both predicates take references. Visiting order is unchanged.
Path: every replayed produce and consume, that is, every block validation.
Risk: \low{}.

\paragraph{0.3 \texttt{Env::put} copies the whole environment on every binding.}
\f{rho-pure-eval/src/env.rs}:31--40.
\begin{lstlisting}
pub fn put(&mut self, a: A) -> Env<A> {
    Env { env_map: { self.env_map.insert(self.level, a);
                     self.env_map.clone() },
          level: self.level + 1, shift: self.shift }
}
\end{lstlisting}
Every caller folds \texttt{put} over a sequence, so binding $k$ values deep-copies
$k(k+1)/2$ \texttt{Par}s. Callers include \texttt{build\_env} on every COMM
(\f{dispatch.rs}:17--19), \texttt{eval\_new} (\f{reduce.rs}:1601--1607), match-case binding
(\f{reduce.rs}:1424--1426) and guard evaluation (\f{matcher/match.rs}:96--98).
\texttt{Env::shift} (line 48) clones the whole map as well, once per receive and once per
substitution under a binder. The signature also has an aliasing hazard: it mutates
\texttt{self} \emph{and} returns a copy, so the caller's environment changes even if the caller
uses only the return value. \emph{Fix, first step:} add \cd{put_owned(mut self, a) -> Self}
and switch the fold sites, which own their accumulator. \emph{Second step:} store the map as
\texttt{Arc<HashMap>} with \texttt{Arc::make\_mut}, or as an \texttt{imbl::HashMap}; then
\texttt{shift} and every per-branch \texttt{env.clone()} become pointer bumps (see 2.1).
Risk: \low{} for the first step (behaviour identical at the fold sites), \medium{} for the
second (type crosses crates).

\paragraph{0.4 \texttt{EventLogIndex::combine} in folds is quadratic in chain length.}
\f{rspace++/.../merger/event_log_index.rs}:443ff. \texttt{combine(x: \&Self, y: \&Self)}
builds the union of all eleven sets with \cd{.union(..).cloned().collect()}. Casper folds
with it (\f{deploy_chain_index.rs}:87, 91; \f{dag_merger.rs}:1408--1435), so an accumulator over
$n$ deploys is re-copied $n$ times. \emph{Fix:} add \cd{combine_into(&mut self, &other)} that
extends in place; keep \texttt{combine} for other callers. Set semantics do not depend on
insertion order. Risk: \low{}--\medium{}.

\paragraph{0.5 \texttt{eval\_match} copies the remaining cases on every failed case.}
\f{reduce.rs}:1513 passes \cd{mat.cases.clone()}; each failed case then does
\cd{state = (_target, case_rem.to_vec())} (lines 1453, 1479). A match with $n$ cases that
falls through to the last one copies $O(n^2)$ case bodies. \emph{Fix:} loop over an index into
\cd{&mat.cases}. Risk: \low{}.

\paragraph{0.6 \texttt{SortedParMap::extend} re-sorts the whole map once per element.}
\f{models/src/rust/sorted_par_map.rs}:43--53. \texttt{insert} clones the map and rebuilds the
sorted form; \texttt{extend} calls \texttt{insert} in a loop, discarding each result, then
rebuilds once more. \texttt{SortedParHashSet} has the same shape. The reducer uses these for the
Rholang \texttt{++}, \texttt{+}, \texttt{-} operations on maps and sets. \emph{Fix:} insert all
pairs, then build the sorted form once. Risk: \medium{}; this is canonical-form code, and the
fix needs a differential test that compares sorted outputs (Section~\ref{sec:guard}).

\paragraph{0.7 Light-finalised-state sync state is copied per block.}
\f{casper/.../engine/lfs_block_requester.rs}:136--148, 339--360. The requester's
\texttt{HashMap} and \texttt{HashSet} state is cloned on every \texttt{add} and \texttt{done},
so sync of $n$ blocks costs $O(n^2)$. This affects only initial sync, so it is recorded here for
completeness and scheduled last. Risk: \low{}.

\subsection{Tier 1: hot path, low risk, mechanical}

Each item below changes ownership only. None changes a value that is hashed, charged, sorted or
logged.

\begin{longtable}{>{\raggedright\arraybackslash}p{0.05\linewidth}>{\raggedright\arraybackslash}p{0.33\linewidth}>{\raggedright\arraybackslash}p{0.52\linewidth}}
\toprule
\# & Location & Change\\
\midrule\endhead
1.1 & \f{reduce.rs}:401--408 (\texttt{eval\_inner} fan-out) &
\cd{terms} is owned and dead after the map; use \cd{into_iter()} instead of
\cd{term.clone()}. Make \texttt{split} take \cd{&Blake2b512Random}: the split is
\emph{required}, the copy of the parent generator is not. Applies to every \texttt{Par} with two or
more terms.\\
1.2 & \f{substitute.rs}:205 & \cd{expr.clone().expr_instance} deep-copies every \texttt{Expr} of
every substituted \texttt{Par} to inspect its variant. Use \cd{match &expr.expr_instance}.
Same pattern at lines 165, 184, 407, 630, 975.\\
1.3 & \f{substitute.rs}:63, 90 & \texttt{substitute\_and\_charge} clones its input only so that
the error path can read \texttt{encoded\_len()}. Add a by-value variant; compute the length
before the move. Charge values and order unchanged.\\
1.4 & \f{reduce.rs}:1725 & \texttt{eval\_expr\_to\_par} deep-copies each expression to match on it.
Match by reference.\\
1.5 & \f{reduce.rs}:1740, 7485 & \texttt{method\_table()} builds a \texttt{HashMap} of about
60 \texttt{String} keys and boxed closures on \emph{every} method call. Replace with a
\texttt{match} on the name (an allocation fix rather than a clone fix).\\
1.6 & \f{reduce.rs}:2038--2916 (32 sites) & \cd{&p1.clone().unwrap()} deep-copies an operand to
borrow it. Use \cd{p1.as_ref().unwrap()}. Similar at 1283, 1442, 1694, 2679.\\
1.7 & \f{reduce.rs}:563, 648--652 & \texttt{produce\_inner} and \texttt{consume\_inner} clone
channel, data, sources, patterns, the whole continuation body and the guard; the originals are
used again only in the \texttt{persistent} branch. Clone only when persistent.\\
1.8 & \f{charging_rspace.rs}:89, 113 & The full RSpace result (continuation body and every matched
datum) is cloned so that \texttt{handle\_result} can read it. Pass by reference.\\
1.9 & \f{dispatch.rs}:53 & \cd{build_env(data_list.clone())}; build the random-state vector first,
then move. The \emph{required} random-state merge order is unchanged.\\
1.10 & \f{space_matcher.rs}:94 & \cd{rollback.push((channel.clone(), original.clone()))} copies
the whole indexed data vector per matched channel. Use \cd{std::mem::replace}.\\
1.11 & \f{space_matcher.rs}:136--140 & Channel and pattern \texttt{Par}s are cloned into pairs for
every candidate continuation probed. Pair references.\\
1.12 & \f{hot_store.rs}:414, 471--474 & Installed (system) continuations are stored by value and
deep-cloned into a fresh \texttt{Arc} on every produce on an installed channel. Store them as
\texttt{Arc}, as the ordinary continuations already are.\\
1.13 & \f{shared/.../key_value_store.rs}:49--52 & \texttt{get\_one} clones the first value
instead of moving it: a second full copy of every value read, including every stored block and
every trie node. Use \cd{values.into_iter().next().flatten()}. Same at
\f{key_value_typed_store_impl.rs}:76.\\
1.14 & \f{interpreter_util.rs}:518--519, 609--610; \f{proto_util.rs}:339 & Replay clones the
full deploy and system-deploy lists once to extract them and again per replay attempt. The
consumer only iterates. Take \cd{&[ProcessedDeploy]}.\\
1.15 & \f{interpreter_util.rs}:386; \f{block_creator.rs}:3482 & \cd{parents.clone()} of
\cd{Vec<BlockMessage>} into \texttt{compute\_parents\_post\_state}; at 386 the original is
dead, and in the proposer it is copied on every retry. Pass a slice.\\
1.16 & \f{block_processor.rs}:542 & \cd{block.to_proto().encode_to_vec().len()} clones the body
twice and encodes it only to record a size metric. Use \texttt{encoded\_len()}, as
\f{block_creator.rs}:3648 already does.\\
1.17 & \f{block_processor.rs}:470, 499, 1097; \f{casper_launch.rs}:302, 322 & Whole blocks cloned
into \texttt{PrettyPrinter::build\_string} for logging. Use the borrowing
\texttt{build\_string\_block\_message(\&block, \ldots)}. No risk.\\
1.18 & \f{conflict_set_merger.rs}:1194 & \texttt{fold\_rejection} clones the balance map for
every branch. Take \cd{&HashMap} and return a new map only on success. Fold order unchanged.\\
1.19 & \f{replay_rspace.rs}:649--653 & \texttt{locked\_consume} deep-copies all data on the
channels only to collect produce hashes, then fetches the same data again. Reuse it.\\
\bottomrule
\end{longtable}

\subsection{Tier 2: structural changes, medium risk, large gains}

\paragraph{2.1 Persistent \texttt{Env}.} After 0.3, every spawned branch in \texttt{eval\_inner},
every \texttt{shift}, and every substitution under a binder still deep-copies the environment.
An \texttt{imbl::HashMap} or \texttt{Arc}-backed map turns all of these into pointer bumps. This
is the largest single structural win in the interpreter. The type is public in
\texttt{rho-pure-eval}, so the change crosses crates.

\paragraph{2.2 \texttt{Arc<Datum>} in the hot store.} \texttt{get\_data}
(\f{hot_store.rs}:631--641) returns the whole \cd{Vec<Datum<A>>} for a channel on every
produce and consume, because \texttt{ShardedMap::get} returns \texttt{.cloned()}
(lines 107--112). On a busy channel this is the largest per-operation copy in RSpace. Add
\texttt{get\_data\_arc} returning \cd{Vec<Arc<Datum<A>>>}, mirroring the existing
\texttt{get\_continuations\_arc}; keep \texttt{get\_data} for other callers. There is a second,
probable effect: after each soft checkpoint the \texttt{imbl} shards are shared with the
snapshot, so the next \texttt{get\_mut} copies a map node, which copies every sibling value in
it. With \texttt{Arc} values those copies become pointer bumps. This second effect should be
confirmed against the \texttt{imbl} source before it is claimed.

\paragraph{2.3 \texttt{Arc<DeployChainIndex>} and \texttt{Arc<BlockIndex>}.} The block-index
cache returns deep copies (\f{interpreter_util.rs}:1317, \cd{(*cached.value()).clone()}), so
every merge deep-copies the index of every block in scope, event logs included. The merger then
clones chains again (\f{dag_merger.rs}:433, 459, 468, 734--737). \texttt{DeployChainIndex} has
hand-written \texttt{Eq}, \texttt{Hash} and \texttt{Ord} over \texttt{deploys\_with\_cost} only,
so \texttt{Arc} wrapping keeps merge order identical by construction. The one mutation,
\texttt{stamp\_prior\_rejections} (\f{dag_merger.rs}:602--604), becomes \texttt{Arc::make\_mut}.

\paragraph{2.4 Borrowing merge accessors.} \f{dag_merger.rs}:1504--1507 defines
\cd{|chain| Ok(chain.state_changes.clone())} and a similar accessor for mergeable channels.
The merger calls them repeatedly per item, branch and rejection option
(\f{conflict_set_merger.rs}:217, 553, 908, 1105). Returning references needs a
\texttt{combine\_ref} on \texttt{StateChange} and a change to the generic \texttt{merge} signature.
\f{dag_merger.rs}:212--213 also clones two growing datum maps per chain; an undo log would
remove that.

\paragraph{2.5 Encode once, send many.} \texttt{stream\_mult}
(\f{grpc_transport_client.rs}:590) streams to each peer separately. Each stream encodes the packet
into the cache, decodes it, clones it, LZ4-compresses the whole block and copies each chunk
(\f{stream_observable.rs}:147, \f{chunker.rs}:22--24). Compress and chunk once, and share the
chunks as \texttt{Bytes} slices. Wire bytes are identical.

\paragraph{2.6 Borrowing protobuf conversion.} \cd{BlockMessage::to_proto(&self)} and
\cd{Body::to_proto} clone every deploy and justification before converting
(\f{casper_message.rs}:834, 956--968, 1508--1514). This runs at store, gossip, hashing and the
size metric. Add \texttt{to\_proto\_ref} methods that map over references. Output bytes are
identical, but because \texttt{hash\_block} uses this path, the change needs a golden-hash test.

\paragraph{2.7 DAG representation without LMDB scans.} Not a clone, but the same kind of waste:
\texttt{get\_representation\_internal}
(\f{block_dag_key_value_storage.rs}:1199--1207) rebuilds \texttt{latest\_messages} and
\texttt{invalid\_blocks} by scanning and decoding LMDB on every call, and
\texttt{admit\_dag\_contains} builds a whole representation to answer one membership query. A
\texttt{contains} method that reads \texttt{dag\_state} is \low{} risk; caching the two
collections in \texttt{DagState} is \medium{} (restart recovery must rebuild them).

\paragraph{2.8 Sorter ownership.} The sort matchers clone whole scored vectors to separate terms
from scores (for example \f{send_sort_matcher.rs}:26,
\cd{sorted_data.clone().into_iter().map(|p| p.term)}), and several clone an
unsorted node only to overwrite the field they just copied. Each recursion level copies
everything below it, which is quadratic in depth. The unzip helper \texttt{split\_scored\_terms}
already exists in \f{par_sort_matcher.rs}. Comparisons and order are untouched, but the code is
the canonical form of every term, so this is graded \medium{} and placed behind the
differential test.

\subsection{Tier 3: not recommended now}

\begin{itemize}
\item \textbf{\texttt{Blake2b256Hash} as \texttt{[u8;32]}.} Removes an allocation per clone across
RSpace, but its derived serde writes a length-prefixed vector. Cold-store keys, history-trie keys
and event-log hashing would change unless a byte-identical custom serialiser is written. \high{}.
\item \textbf{\texttt{Bytes} or \texttt{Arc} inside \texttt{rhoapi}.} Would make \texttt{Par}
clones cheap at the root, but touches the serde and schema derives and the
\texttt{locally\_free} handling in \f{models/build.rs}. \high{}.
\item \textbf{Removing the second sort of maps and sets.} Depends on sort idempotence and on
identical deduplication semantics. \high{} without a proof.
\item \textbf{\texttt{continuation\_identity}} (\f{hot_store.rs}:1012). It formats a
\cd{{:?}} string of every continuation on a channel at each \texttt{put\_continuation}. Not a
clone, but expensive; it decides store contents and therefore the state hash, so a replacement
needs an equivalence argument. \medium{}--\high{}.
\item \textbf{Zero-copy LMDB reads.} Would change the on-disk format. \high{}.
\end{itemize}

\subsection{Dismissed}
Start-up wiring in \f{node/.../setup.rs}, \f{node_runtime.rs} and \f{casper_launch.rs};
\texttt{Arc} handle bumps for spawned tasks; clones of \texttt{BlockHash}, \texttt{StateHash} and
\texttt{Validator}; per-request clones in \f{web_api.rs} (small strings and hashes); and clique
oracle map clones (small validator sets). These account for well over a third of the census and
for a negligible part of the cost.

%=====================================================================
\section{Microbenchmarks of the patterns}
\label{sec:bench}

Four patterns were reproduced in a standalone Rust program with the same shape as the
repository code, and timed against the proposed fix (release build, stable toolchain, single
thread). A synthetic \texttt{Tree} with nested \texttt{Vec}s stands in for \texttt{Par}. The
numbers measure the pattern, not the node; they show growth rates, not production latency.

\begin{table}[H]
\centering\small
\begin{tabular}{llrrr}
\toprule
Pattern & Size & Current & Fixed & Ratio\\
\midrule
Stream reassembly (0.1) & 1 MiB  & 1.63 ms & 0.035 ms & 46$\times$\\
                        & 10 MiB & 84.5 ms & 0.82 ms & 103$\times$\\
                        & 64 MiB & 5{,}562 ms & 43.9 ms & 127$\times$\\
\midrule
\texttt{Env::put} fold (0.3), 85-node terms & 2 bindings & 14.0\,\textmu s & 5.1\,\textmu s & 2.7$\times$\\
                       & 8 bindings & 126.7\,\textmu s & 23.1\,\textmu s & 5.5$\times$\\
                       & 32 bindings & 2{,}160\,\textmu s & 128\,\textmu s & 16.9$\times$\\
\midrule
Replay COMM scan (0.2), match last & 8 COMMs & 5.3\,\textmu s & 0.6\,\textmu s & 9.5$\times$\\
                       & 64 COMMs & 542\,\textmu s & 8.1\,\textmu s & 67$\times$\\
                       & 256 COMMs & 7{,}831\,\textmu s & 28.7\,\textmu s & 273$\times$\\
\midrule
Fan-out, 16 subterms (1.1) & 121-node terms & 134.7\,\textmu s & 62.6\,\textmu s & 2.2$\times$\\
\bottomrule
\end{tabular}
\caption{Pattern microbenchmarks. ``Fixed'' is the change proposed in this note.}
\end{table}

The first three rows grow with input size, as the analysis predicts; the last is a constant
factor, as expected for a single redundant copy. The stream result deserves emphasis: at 64\,MiB
the current pattern spends more than five seconds copying buffers that the fixed pattern
appends in place.

%=====================================================================
\section{Invariants that must not move}
\label{sec:guard}

Some clones in the hot paths are \emph{required}, and the plan must leave them alone.

\begin{enumerate}
\item \textbf{Random-number splitting.} Each parallel branch receives a generator split from its
parent by index (\f{reduce.rs}:408). The split determines every unforgeable name and must stay.
Only the extra copy of the parent before splitting goes.
\item \textbf{Cost accounting.} The amount and order of every charge must stay identical. Items
1.3, 1.8 and the cost helpers that take a whole value to read its length change only how a length
is obtained.
\item \textbf{Canonical form.} Sorting defines term equality and hashes. Changes to the sorter
(0.6, 2.8) must produce byte-identical output.
\item \textbf{Hash and serialisation bytes.} \texttt{hash\_block}, event-log hashes, trie keys and
LMDB values must not change.
\item \textbf{Replay semantics.} COMM matching order in replay must be preserved; 0.2 keeps the
iteration order.
\end{enumerate}

\paragraph{Guards.} Three tests make the medium-risk work safe.
\begin{itemize}
\item \emph{Golden state hashes.} Run a fixed corpus of Rholang deploys (the existing test
contracts, \texttt{PoS.rhox}, and the registry) from genesis and record the post-state hash and
event-log hash of each. Every pull request in Tiers 1 and 2 must reproduce them exactly.
\item \emph{Proto round trip.} For a corpus of blocks, assert that new and old
\texttt{to\_proto} produce identical encoded bytes and identical block hashes.
\item \emph{Sorter differential.} For random \texttt{Par}s (the \texttt{fuzz/} crate already has
generators), assert that old and new sort matchers produce identical terms and score trees.
\end{itemize}

%=====================================================================
\section{Proposed sequence}
\label{sec:measure}

Each step is one small pull request to \texttt{dev}, with the \texttt{ci-heavy} label, so the
integration pipeline runs before the merge queue.

\begin{enumerate}
\item \textbf{Guards first.} Land the golden-hash and proto round-trip tests, and record a
baseline with \f{scripts/bench/casper-soak.sh} and \f{latency-benchmark.sh}: block validation
time, replay time, propose time, peak resident memory, and the
\cd{hot-store.get_data.calls} and \cd{hot-store.get_data.history_fill} metrics.
\item \textbf{Asymptotic, low risk:} 0.1 (stream), 0.2 (replay scan), 0.3 first step
(\texttt{put\_owned}), 0.5 (match cases). Four independent PRs.
\item \textbf{Interpreter mechanical:} 1.1--1.9. Mostly within \f{reduce.rs},
\f{substitute.rs}, \f{dispatch.rs}, \f{charging_rspace.rs}.
\item \textbf{RSpace and storage mechanical:} 1.10--1.13, 1.19.
\item \textbf{Casper mechanical:} 1.14--1.18. Together with 0.1, these remove about six of the
roughly eight full body copies an incoming block undergoes between receipt and storage.
\item \textbf{Structural:} 0.4, 2.1, 2.2, 2.3, 2.4 --- one at a time, each with a fresh soak
comparison.
\item \textbf{Transport and model:} 2.5, 2.6, 2.7, then 0.6 and 2.8 behind the sorter
differential.
\item \textbf{Prevention.} Add an advisory CI job that runs Clippy with
\cd{-W clippy::redundant_clone} and \cd{-W clippy::needless_pass_by_value} on the interpreter, RSpace and
merge modules. The current \texttt{cargo clippy} step uses default lints, and
\texttt{redundant\_clone} is a nursery lint, so it does not run today. Keep the job
non-blocking at first; the pedantic lint produces noise on cold code.
\end{enumerate}

\paragraph{Expected effect.} Without profiles, an honest estimate is qualitative. Steps 2--5 are
mechanical and should remove most of the avoidable deep copies on the evaluation, replay and
block-intake paths. The largest wall-clock gains are expected where growth rates change: large
block transfer (0.1), replay of blocks with many COMMs on a channel (0.2), contracts that bind
many names or match on many cases (0.3, 0.5), and merges over long deploy chains (0.4, 2.3). A
soak comparison after step 5 will show whether Tier~2 is needed for current workloads or can wait.

%=====================================================================
\section{Summary}

The clone count is high, but most clones are cheap: the consensus hash types are already
reference-counted, and the DAG representation is already persistent. The cost is concentrated
where the Scala port copies deep values to imitate immutable-collection semantics. Seven sites
are asymptotic defects, and four of them are low-risk, few-line fixes. About twenty further hot-path
sites need only a change of ownership. The structural changes --- a persistent \texttt{Env},
\texttt{Arc}-shared data in the hot store and in the merge indices --- extend a pattern the
code base has already applied successfully to continuations. All of this can proceed without
touching random-number splitting, cost order, canonical sorting, hash bytes or replay order,
provided the golden-hash, round-trip and sorter-differential guards land first.

\appendix
\section{Reference table of cited locations}
\small
\begin{longtable}{>{\raggedright\arraybackslash}p{0.08\linewidth}>{\raggedright\arraybackslash}p{0.55\linewidth}>{\raggedright\arraybackslash}p{0.12\linewidth}>{\raggedright\arraybackslash}p{0.15\linewidth}}
\toprule
Item & File : line & Path & Risk\\
\midrule\endhead
0.1 & \f{comm/src/rust/transport/stream_handler.rs}:482--491, 393 & per block & \low\\
0.2 & \f{rspace++/src/rspace/replay_rspace.rs}:1126--1153, 966, 973 & per replay op & \low\\
0.3 & \f{rho-pure-eval/src/env.rs}:31--51 & per COMM & \low/\medium\\
0.4 & \f{rspace++/src/rspace/merger/event_log_index.rs}:443ff & per merge & \low/\medium\\
0.5 & \f{rholang/src/rust/interpreter/reduce.rs}:1453, 1479, 1513 & per match & \low\\
0.6 & \f{models/src/rust/sorted_par_map.rs}:43--69 & map/set ops & \medium\\
0.7 & \f{casper/src/rust/engine/lfs_block_requester.rs}:136--148 & sync only & \low\\
1.1--1.9 & \f{reduce.rs}, \f{substitute.rs}, \f{dispatch.rs}, \f{charging_rspace.rs} & per eval / COMM & \low\\
1.10--1.13 & \f{space_matcher.rs}, \f{hot_store.rs}, \f{key_value_store.rs} & per op / read & \low\\
1.14--1.18 & \f{interpreter_util.rs}, \f{block_creator.rs}, \f{block_processor.rs}, \f{conflict_set_merger.rs} & per block & \low\\
2.1 & \f{rho-pure-eval/src/env.rs} & per eval & \medium\\
2.2 & \f{rspace++/src/rspace/hot_store.rs}:107--112, 631--641 & per op & \medium\\
2.3 & \f{interpreter_util.rs}:1317; \f{dag_merger.rs}:433--737 & per merge & \medium\\
2.4 & \f{dag_merger.rs}:212, 1504--1507 & per merge & \medium\\
2.5 & \f{grpc_transport_client.rs}:590; \f{chunker.rs}:22 & per broadcast & \low/\medium\\
2.6 & \f{models/src/rust/casper/protocol/casper_message.rs}:834--1514 & per block & \medium\\
2.7 & \f{block_dag_key_value_storage.rs}:1199--1207 & per block & \low/\medium\\
2.8 & \f{models/src/rust/rholang/sorter/*} & per substitution & \medium\\
\bottomrule
\end{longtable}

\end{document}
